\chapter*{Abstract}
\addcontentsline{toc}{chapter}{Abstract}

Reconstructing a multi-stage cyber attack from host audit logs is still largely manual work. Structured prompting topologies (Chain-, Tree- and Graph-of-Thoughts) improve language-model reasoning on short, well-defined tasks, but it is unknown whether the extra structure helps when the evidence is a long and noisy provenance log. This thesis isolates that question. We build a pipeline in which ground truth, evidence construction, reduction, serialization, prompt, parser and scorer are identical for the three topologies, so the topology is the only variable. It is evaluated on five attack scenarios from the NODLINK Simulated-Data release (89 annotated attack events) with locally hosted open models served by vLLM and no training or fine-tuning. Our hypothesis was that the Aggregate operation of Graph-of-Thoughts would raise stage-level recall.

On the metric named in the hypothesis, stage recall, GoT was not reliably better than CoT or ToT: 2 of 15 paired comparisons fell below $p=0.05$ before correction and none survived it. GoT did recover more events (uncorrected $p<0.05$ in 5 of 10 comparisons) at lower precision for the smaller models. Reading raw outputs against ground truth explained why: only 15.3\% of true attack events were ever detected, and exact stage labels were right about 20\% of the time, because each model carries a stable stage-label prior that does not match the data. Adding MITRE ATT\&CK stage definitions to the shared prompt doubled exact stage recall (0.056 to 0.123, Wilcoxon $p=0.0015$); per-process triage reduced false positives (precision 0.245 to 0.329, $p=0.04$); and scaling Qwen 2.5 from 7B to 32B raised exact stage recall from 0.15 to 0.30 ($p=0.012$). Graph-of-Thoughts had the highest event $F_1$ and stage recall for both large models. A security-tuned 8B model scored below the general model (event $F_1$ 0.101 against 0.295), partly because 38\% of its runs gave no parseable answer, and did not separate the topologies. On the longest inputs a 70B model scored below an 8B model, which shows input length is a confound in topology comparisons. We conclude that at the 7--70B scale, detection recall and stage calibration matter more than reasoning topology, and we propose a segmented-Chain-of-Thought control to separate the benefit of splitting from the benefit of aggregating.
